% output=pdftex

\setupbodyfont[loc,ppl]

\startbuffer[font]
\startMPenvironment
  \setupbodyfont[loc,ppl]
\stopMPenvironment
\stopbuffer

\getbuffer[font]

\def\compoundhyphen{-}

\setuplayout
  [backspace=1.5cm,
   topspace=1.5cm,
   width=middle,
   height=middle]

\setupwhitespace
  [big]

\setupcolors
  [state=start]

\setuptyping
  [color=red]

\setupheadertexts[makempy]
\setupfootertexts[pagenumber]

\usemodule[chart,abr-01]

\startFLOWchart[makempy]
  \setupFLOWcharts
    [breedte=6\bodyfontsize,
     hoogte=3\bodyfontsize,
     dx=\bodyfontsize,
     dy=\bodyfontsize]
  \setupFLOWshapes
    [framecolor=blue]
  \setupFLOWlines
    [color=red]
  \startFLOWcell
    \name       {metapost}
    \location   {1,1}
    \shape      {action}
    \text       {\hbox{\MetaPost}}
    \connection [bt] {mpo file}
  \stopFLOWcell
  \startFLOWcell
    \name       {mpo file}
    \location   {1,2}
    \shape      {action}
    \text       {\sl mpo file}
    \connection [bt] {makempy}
  \stopFLOWcell
  \startFLOWcell
    \name       {makempy}
    \location   {1,3}
    \shape      {action}
    \text       {makempy}
    \connection [bt] {mpy file}
    \connection [rl] {pdftex}
  \stopFLOWcell
  \startFLOWcell
    \name       {mpy file}
    \location   {1,4}
    \shape      {action}
    \text       {\sl mpy file}
    \connection [bl] {metapost}
  \stopFLOWcell
  \startFLOWcell
    \name       {pdftex}
    \location   {2,3}
    \shape      {action}
    \text       {pdf\TeX}
    \connection [rl] {pdftops}
  \stopFLOWcell
  \startFLOWcell
    \name       {pdftops}
    \location   {3,3}
    \shape      {action}
    \text       {pdftops}
    \connection [rl] {pstoedit}
  \stopFLOWcell
  \startFLOWcell
    \name       {pstoedit}
    \location   {4,3}
    \shape      {action}
    \text       {pstoedit}
    \connection [rpb] {makempy}
  \stopFLOWcell
  \startFLOWcell
    \name       {tex}
    \location   {2,1}
    \shape      {action}
    \text       {\TeX}
    \connection [bt] {pdftex}
  \stopFLOWcell
  \startFLOWcell
    \name       {context}
    \location   {3,1}
    \shape      {action}
    \text       {\ConTeXt}
    \connection [bt] {pdftex}
  \stopFLOWcell
  \startFLOWcell
    \name       {latex}
    \location   {4,1}
    \shape      {action}
    \text       {\LaTeX}
    \connection [bt] {pdftex}
  \stopFLOWcell
\stopFLOWchart

\definecolor[gray] [s=.7]
\definecolor[red]  [r=.7]
\definecolor[green][g=.7]
\definecolor[blue] [b=.7]

\startbuffer[head]
\startuseMPgraphic{text}
  graphictext \MPstring{string} scaled 4
    withdrawcolor \MPcolor{blue} withfillcolor \MPcolor{gray}
    withpen pencircle scaled 2pt ;
\stopuseMPgraphic

\startuseMPgraphic{number}
  graphictext \MPstring{string} scaled 10
    withdrawcolor \MPcolor{red} withfillcolor \MPcolor{gray}
    withpen pencircle scaled 2pt ;
\stopuseMPgraphic

\def\TextCommand#1%
  {\setMPtext{string}{#1}%
   \framed
     [frame=off,offset=overlay,height=3cm,width=\hsize]
     {\useMPgraphic{text}}}

\def\NumberCommand#1%
  {\expanded{\setMPtext{string}{\currentheadnumber}}%
   \hsmash
     {\framed
        [frame=off,offset=overlay,height=3cm,width=\hsize]
        {\useMPgraphic{number}}}}

\setuphead
  [section]
  [distance=0pt,
   textcommand=\TextCommand,
   numbercommand=\NumberCommand]
\stopbuffer

\getbuffer[head]

\starttext

% \startuseMPgraphic{front make}
%   graphictext
%     "\bf MAKE"
%     scaled 8
%     zscaled (1,2)
%     withdrawcolor \MPcolor{gray} % \MPcolor{blue}
%     withfillcolor \MPcolor{blue} % \MPcolor{gray}
%     withpen pencircle scaled 10pt ;
% \stopuseMPgraphic
% 
% \startuseMPgraphic{front mpy}
%   graphictext
%     "\bf MPY"
%     scaled 5
%     zscaled (1,2)
%     withdrawcolor \MPcolor{gray} % \MPcolor{red}
%     withfillcolor \MPcolor{red}  % \MPcolor{gray}
%     withpen pencircle scaled 6pt ;
% \stopuseMPgraphic
% 
% \startuseMPgraphic{hans}
%   graphictext
%     "Hans"
%     scaled 6
%     withdrawcolor \MPcolor{gray} % \MPcolor{red}
%     withfillcolor \MPcolor{red}  % \MPcolor{gray}
%     withpen pencircle scaled 4pt ;
% \stopuseMPgraphic
% 
% \startuseMPgraphic{hagen}
%   graphictext
%     "Hagen"
%     scaled 6
%     withdrawcolor \MPcolor{gray} % \MPcolor{red}
%     withfillcolor \MPcolor{red}  % \MPcolor{gray}
%     withpen pencircle scaled 4pt ;
% \stopuseMPgraphic
% 
% \startuseMPgraphic{page}
%   StartPage ;
%     fill Page
%       withcolor .5white ;
%     draw Page enlarged -12pt
%       withpen pencircle scaled 5pt
%       withcolor \MPcolor{gray} ;
%   StopPage ;
% \stopuseMPgraphic
% 
% \defineoverlay[page] [\useMPgraphic{page}]
% \defineoverlay[make] [\useMPgraphic{front make}]
% 
% \setupbackgrounds[page][background={page,make}]
% 
% \startstandardmakeup
%   \hbox to \hsize{\useMPgraphic{front mpy}\hss}
%   \vfill
%   \hbox to \hsize{\hss\useMPgraphic{hans}}
%   \vskip24pt
%   \hbox to \hsize{\hss\useMPgraphic{hagen}}
%   \vskip-24pt
% \stopstandardmakeup
% 
% \setupbackgrounds[page][background=]

\startuseMPgraphic{metapost}
  graphictext
    "\bf MetaPost"
    scaled 8
    withdrawcolor .5white 
    withfillcolor .8blue 
    withpen pencircle scaled 10pt ;
\stopuseMPgraphic

\startuseMPgraphic{creating}
  graphictext
    "\bf creating"
    scaled 5
    withdrawcolor .5white 
    withfillcolor .8red 
    withpen pencircle scaled 6pt ;
\stopuseMPgraphic

\startuseMPgraphic{outlines}
  graphictext
    "\bf outlines"
    scaled 5
    withdrawcolor .5white 
    withfillcolor .8red 
    withpen pencircle scaled 6pt ;
\stopuseMPgraphic

\startuseMPgraphic{makempy}
  graphictext
    "\bf makempy manual"
    scaled 4
    withdrawcolor .5white 
    withfillcolor .8blue 
    withpen pencircle scaled 4pt ;
\stopuseMPgraphic

\startuseMPgraphic{hans}
  graphictext
    "\bf\strut Hans Hagen"
    scaled 4
    withdrawcolor .5white 
    withfillcolor .8red
    withpen pencircle scaled 4pt ;
\stopuseMPgraphic

\startuseMPgraphic{page}
  StartPage ;
    numeric u ; 
    pair p ; 
    fill Page
      withcolor .5white ;
    for i=1 upto 512 : 
      p := center Page randomized (PaperWidth-75pt,PaperHeight-75pt) ;
      u := 10pt + uniformdeviate 40pt ; 
      draw p withcolor .8white withpen pencircle scaled u ;
      draw p withcolor yellow randomized (.5,.8) withpen pencircle scaled (u-5pt) ;
    endfor ; 
    draw Page enlarged -12pt
      withpen pencircle scaled 5pt
      withcolor .8white ;
  StopPage ;
\stopuseMPgraphic

\defineoverlay[page] [\useMPgraphic{page}]

\setupbackgrounds[page][background=page]

\startstandardmakeup
  \hbox to \hsize{\hss\useMPgraphic{creating}}
  \hbox to \hsize{\hss\useMPgraphic{metapost}}
  \vskip6pt
  \hbox to \hsize{\hss\useMPgraphic{outlines}}
  \vfill 
  \hbox to \hsize{\hss\useMPgraphic{makempy}}
  \vskip6pt
  \hbox to \hsize{\hss\useMPgraphic{hans}}
\stopstandardmakeup

\setupbackgrounds[page][background=]

\section{introduction}

You can use \METAPOST\ to create graphics in a declarative
manner. Although there are tools to create \METAPOST\
graphics interactively, by nature the input is a script
defining the graphic.

Plain \METAPOST\ can handle text directly, as strings, or as
pictures distilled from \TEX\ output. In both cases, the
text is composed of individual characters, and these somehow
will end up in the \POSTSCRIPT\ file generated by \METAPOST.

Although one should not see \METAPOST\ as a full blown
system for doing fancy graphic, its usage, especially in a
dynamic document as created by \TEX, sometimes demands a
more sophisticated way of handling text. The section titles
in this document are an example of this.

In this document I will discuss a way to import text (as
typeset by \TEX) into a \METAPOST\ graphic. Although
primarily written for usage with \CONTEXT, this method is
quite generic and also works well with plain \TEX,
\LATEX\ or others.

\section{method}

Including text as graphics is far from trivial. First it
has to be typeset, and of course we want to use \TEX\ for
that, and in our case \PDFTEX\ suits well. Next we need to
convert the typeset text to graphics, or actually outlines.
For that step we will use Wolfgang Glunz's \type
{pstoedit}, which can produce \METAPOST\ code from
\POSTSCRIPT\ code. This program falls back on \GHOSTSCRIPT\
for creating the curves out of the fonts. Because (at least
currently) handling \PDF\ directly is not working as
expected, we convert the \PDF\ file to \POSTSCRIPT\ with
Derek Noonburg's \type {pdftops} program. This process is
illustrated in \in {figure} [fig:process].

\placefigure
  [here]
  [fig:process]
  {The process of conversion.}
  {\FLOWchart[makempy]}

\section{tools}

During a \METAPOST\ run, the text to be processed is
written to a file with the suffix \type {mpo}. This file,
which only has \TEX\ directives, is converted to a file
with suffix \type {mpy} which holds \METAPOST\ code. The
conversion is taken care of by a \PERL\ script \type
{makempy}. This script generates some intermediate files
that are fed into the programs mentioned.

In a next \METAPOST\ run, the \type {mpy} file is read in
each time a graphic is needed, and the curves are processed
in a special way that permits us to fill and|/|draw them.
We can apply transformations and use colors and pens. The
macros that take care of this are collected in the
\METAPOST\ file \type {mp-grph.mp}, which is part of the
\CONTEXT\ distribution, but in itself is an independent
part of the \METAFUN\ suite.

\section{usage}

We will now demonstrate this mechanism using a typical \TEX\
example.

\startbuffer[einstein]
input mp-grph ;

beginfig(1) ;
  graphictext
    "$e=mc^2$"
    scaled 10
    dashed evenly
    withdrawcolor .7blue
    withfillcolor .7white
    withpen pencircle scaled 2pt ;
endfig ;
\stopbuffer

\typebuffer[einstein]

This example introduces a \METAPOST\ macro \type
{graphictext}. The first argument of this macro is a
string. After this string you can specify transform
operations like a scaling factor, a shift, a rotation or a
combination of these. In addition to the normal color, pen
and dash specifications, there are two special color
operators. When \type {fillcolor} is provided, the graphic
will be also be filled, otherwise it will be an outline.

\startlinecorrection
\hbox to \hsize \bgroup \hss \startMPcode
  graphictext
    "$e=mc^{\scriptscriptstyle2}$" % nicer
    scaled 10
    dashed evenly
    withdrawcolor .7blue
    withfillcolor .7white
    withpen pencircle scaled 2pt ;
\stopMPcode \hss \egroup
\stoplinecorrection

In this example, the string to be processed by \TEX\ is given
between \type {""}. When using \METAPOST\ this way, you
need to be aware of the fact that there can be no line feeds
in a string, so a more complicated formula has to be specified
as follows:

\starttyping
graphictext
  ("$$\pmatrix{D_x&D_y&1\cr} =" &
   "  \pmatrix{U_x&U_y&1\cr}  " &
   "  \pmatrix{s_x&r_x&0\cr   " &
   "           r_y&s_y&0\cr   " &
   "           t_x&t_y&1\cr}$$")
  ....
\stoptyping

So, we split the string into lines, which we paste with
\type {&} and surround by \type {()}. However, when in
\CONTEXT\ you want to include such more complicated \TEX\
code in a graphic that is defined in the source, you need
to be aware of unwanted expansion. This is why in \CONTEXT\
we can best define the string separately, which also keeps
the graphic more readable. There we have:

\startbuffer[one]
\setMPtext{formula}
  {$$\pmatrix{D_x&D_y&1\cr} =
     \pmatrix{U_x&U_y&1\cr}
     \pmatrix{s_x&r_x&0\cr
              r_y&s_y&0\cr
              t_x&t_y&1\cr}$$}
\stopbuffer

\typebuffer[one]

This string can be used as follows:

\startbuffer[two]
\startMPcode
  graphictext
    \MPstring{formula}
    scaled 2
    withfillcolor white
    withdrawcolor .7blue
    withpen pencircle scaled .75pt ;
\stopMPcode
\stopbuffer

\typebuffer[two]

As you can see from the graphic below, you need to adapt the
size of the pen to the scale.

\startlinecorrection
\hbox to \hsize{\hss\getbuffer[one,two]\hss}
\stoplinecorrection

In some cases a glyph is composed of multiple outlines.
Therefore, by default, the fill is applied after the draw.
As a result, the perceived pen is smaller than the specified
one. One can reverse drawing and filling by supplying the
keyword \type {reversefill} as demonstrated in the next
example.

\startbuffer[two]
\startMPcode
  graphictext
    \MPstring{formula}
    scaled 2
    reversefill
    withcolor .7blue
    withpen pencircle scaled .25pt ;
\stopMPcode
\stopbuffer

\typebuffer[two]

As you can see clearly here, big delimiters are composed of
pieces.

\startlinecorrection
\hbox to \hsize{\hss\getbuffer[one,two]\hss}
\stoplinecorrection

\startbuffer[a]
\startMPcode
  graphictext "MP" scaled 8
    withdrawcolor .7blue withfillcolor .7white
    dashed evenly withpen pencircle scaled 2pt ;
\stopMPcode
\stopbuffer

\startbuffer[b]
\startMPcode
  graphictext "MP" scaled 8 reversefill
    withdrawcolor .7blue withfillcolor .7white
    dashed evenly withpen pencircle scaled 1pt ;
\stopMPcode
\stopbuffer

\startbuffer[c]
\startMPcode
  graphictext "MP" scaled 8 outlinefill
    withdrawcolor .7blue withfillcolor .7white
    dashed evenly withpen pencircle scaled 2pt ;
\stopMPcode
\stopbuffer

In addition to the default method and the reversed fill, we
also have an outline fill. The differences show up clearly
if we use a dashed line:

\typebuffer[a]

In the reversed fill, we will get a thinner line, since the
fill covers half the line.

\typebuffer[b]

The outline fill method will extend the fill to the line,
which sometimes gives nicer results with dashed shapes.

\typebuffer[c]

These three definitions result in the following graphics.
The third alternative is the least efficient because the
paths are duplicated.

\startlinecorrection
\hbox to \hsize \bgroup \hss \startcombination[3*1]
  {\getbuffer[one,a]} {default fill}
  {\getbuffer[one,b]} {reverse fill}
  {\getbuffer[one,c]} {outline fill}
\stopcombination \hss \egroup
\stoplinecorrection

In the first example of this section, we use plain \TEX,
while the second and following examples depend on \CONTEXT.
You can specify a \TEX\ back||end using the following
directive:

\starttyping
graphictextformat := "latex" ;
\stoptyping

In \CONTEXT\ you can specify environments in the normal way
\footnote {For more information on what is normal, you can
consult the \METAFUN\ manual.} like:

\starttyping
\startMPenvironment
\setupbodyfont[pos,11pt]
\stopMPenvironment
\stoptyping

In a stand||alone \METAPOST\ graphic, you can pass
information to \TEX\ using \type {graphictextdirective}:

\starttyping
graphictextformat := "latex" ;
graphixtextdirective "\documentclass[]{article}" ;
\stoptyping

You don't have to provide document structuring commands,
since these are added by \type {makempy}.

You can also use \type {pstoedit} independently, in which case
each page becomes a figure. You can load such a figure using
the \type {loadfigure} macro, as in:

\starttyping
loadfigure("filename.suffix", 4) ;
\stoptyping

This macro makes sure that the content of figure~4, when
present, is added to \type {currentpicture}. You can also
use this macro to include figures defined in other
\METAPOST\ files, so that you can build libraries.

\section{switches}

The \PERL\ script \type {makempy} uses the following
programs:

\starttabulate
\NC gs       \NC a general purpose \POSTSCRIPT\ processor \NC \NR
\NC pdftops  \NC a \PDF\ to \POSTSCRIPT\ converter from
                 the xpdf suite                           \NC \NR
\NC pstoedit \NC a \POSTSCRIPT\ to whatever vector format
                 converter (often comes with \GHOSTVIEW)  \NC \NR
\stoptabulate

When set, the environment variables \type {GS} and \type
{GS_PROG} will be used to identify \GHOSTSCRIPT. Of course
you need to have \PERL\ on your system to run \type
{makempy} as well as \PDFTEX\ (which nowadays is part of
most \TEX\ distributions). If needed, the names of the
other programs that are used can be set with \type
{PDFTOPS}, \type {ACROREAD} and \type {PSTOEDIT}. 

The \type {makempy} script does its work without user
intervention. When it deduces that there are no changes, it
will not run at all, otherwise it will report its steps
and|/|or encountered error. The script takes one argument:

\starttyping
makempy filename
\stoptyping

A suffix, when provided, is stripped. The following switches
are recognized:

\starttabulate[|Tl|p|]
\NC --help        \NC returns some basic help information        \NC \NR
\NC --silent      \NC don't report status messages               \NC \NR
\NC --force       \NC always process the file (no checksum test) \NC \NR
\NC --noclean     \NC don't remove temporary files when finished \NC \NR
\NC --acrobat     \NC use \ACROBAT\ for conversion (only unix)   \NC \NR
\NC --pdftops     \NC use \PDFTOPS\  for conversion              \NC \NR
\NC --ghostscript \NC use \GHOSTSCRIPT\ for conversion           \NC \NR
\stoptabulate

The last few switches are handy when you encounter errors or
get unexpected results. In that case you can do as follows:

\starttyping
mpost   yourfile
makempy yourfile --force --noclean
\stoptyping

after that you can take a look at the files whose name
start with \type {mpy-yourfile}.

\starttyping
mpy-yourfile.tex | .pdf | .log | .pos | .tmp
\stoptyping

When using \CONTEXT, keep in mind that \TEXEXEC\ uses files
with names as \type {mpgraph} which can be prefixed by \type
{yourfile}, so there we get:

\starttyping
mpy-yourfile-mpgraph.tex | .pdf | .log | .pos | .tmp
\stoptyping

At this moment we always use \PDFTEX\ combined with \type
{pdftops} since it guarantees that we get everything we need
in the file. Experiments with \type {dvips} and \type
{dvipsone} were unsatisfying with respect to page
dimensions and font conversion.

\section{remarks}

In order to achieve optimal results, a few tricks are used,
some of which don't deserve a beauty price.

Normally a typeset text will use font sizes between 10 and
30 points, because that's the way \TEX\ systems are set up.
When such small shapes as font glyphs are converted to
curves, we loose quite some accuracy and the results will be
quite awful. For this reason, behind the screens, we scale the
graphic up and down by a factor 10.

Because \type {pdftops} has problems with non standard paper
sizes, we specify a pretty large working area when
converting the \PDF\ file to \POSTSCRIPT. The \type
{pstoedit} program generates tight code, so there we will
loose unwanted white space.

\section{example}

In order to demonstrate how well \TEX\ and \METAPOST\ can
work together, I will give an example that uses a typeset
paragraph. This example uses \CONTEXT, but you can of
course make an independent \METAPOST\ figure instead. First
we define a graphic:

\startbuffer
\startuseMPgraphic{quotation}
  graphictext
    \MPstring{text}
    scaled 2
    withdrawcolor .7blue
    withfillcolor .7white
    withpen pencircle scaled 1pt ;
  picture one ; one := currentpicture ; currentpicture := nullpicture ;
  graphictext
    \MPstring{author}
    scaled 4
    withdrawcolor .7red
    withfillcolor .7white
    withpen pencircle scaled 2pt ;
  picture two ; two := currentpicture ; currentpicture := nullpicture ;
  currentpicture := one ;
  addto currentpicture also two
    shifted  lrcorner one
    shifted - 1.125 lrcorner two
    shifted (0, - 1.250 * ypart urcorner two)  ;
  setbounds currentpicture to boundingbox currentpicture enlarged 12pt ;
\stopuseMPgraphic
\stopbuffer

\typebuffer \getbuffer

In this graphic, we have two text fragments, the first one is
a text, the second one the name of the author. We link the
quotation and author into this graphic using the following
definitions:

\startbuffer
\setMPtext{text}  {\vbox{\hsize 8.5cm \input zapf }}
\setMPtext{author}{\hbox{\sl Hermann Zapf}}
\stopbuffer

\typebuffer \getbuffer

These definitions assume that the file \type {zapf.tex} is
present on the system. The graphic, which is shown on the
last page, can now be typeset using the following call:

\startbuffer
\useMPgraphic{quotation}
\stopbuffer

\typebuffer

For those not familiar with \CONTEXT: the graphic data is
written to a file, and processed either directly or between
\TEX\ runs. If you want to know more about the way \CONTEXT\
can interact with \METAPOST, you may want to take a look at
the \METAFUN\ manual (which also discusses \METAPOST\ in
detail).

Of course you can also process this quotation as stand alone
graphic, in which case the \METAPOST\ code goes between
\type {beginfig}||\type {endfig}.

For \CONTEXT\ users who want to know how the section headers
in this document are defined, I let their definition follow
here.

\typebuffer[head]

There are two points worth noticing. We define two hooks
into the section header command. In the macro that deals
with the section number, we don't use the number as
provided by the argument \type {#1}, but use the raw number
instead. We do so, because the content of this argument
contains label related macros, which don't operate that
well when processed independently from this document. We
need to expand the number, because the graphic is processed
in an independent run.

The second trick concerns the \type {hsmash}. Because the
number and text are placed after each other, we explicitly
have to set the distance to zero, as well as makes sure that
the number does not consume space. By using the \type
{\framed} macro, we get both text and number nicely
centered.

We use default font size but pass the colors (red, blue and
gray) from this particular document. The bodyfont
environment is passed on to \METAPOST's \TEX\ run by saying:

\typebuffer[font]

\section{resources}

The \PERL\ script \type {makempy} as well as the \METAPOST\
module \type {mp-grph.mp} are part of the \CONTEXT\
distribution. You can find the latest version in the zipped
archive that can be fetched from the download page at:

\starttyping
www.pragma-ade.com
\stoptyping

or one of the mirrors. You can also find \CONTEXT\ in the
main stream \TEX\ distributions. If you encounter problems,
you may contact:

\startlines
Hans Hagen
PRAGMA-ADE, Hasselt, NL
j.hagen@pragma-ade.com
\stoplines

But for support you can best go to the \CONTEXT\ mailing
list:

\starttyping
ntg-context@ntg.nl
\stoptyping

This manual is typeset at \currentdate.

\page

\startlinecorrection
\getbuffer
\stoplinecorrection

I sincerely hope that Hermann Zapf will forgive me for
mis||using his quotation |<|from an article on {\sl
HZ}||optimization|>| as well as for manipulating his
Palatino font in this way.

\stoptext
